\documentclass[11pt,a4paper]{article}

% --- Encoding ---
\usepackage[utf8]{inputenc}
\usepackage[T1]{fontenc}
\usepackage[spanish,es-noshorthands]{babel}
\usepackage{lmodern}

% --- Layout ---
\usepackage[margin=2.5cm]{geometry}
\usepackage{parskip}
\setlength{\parindent}{0pt}
\setlength{\parskip}{0.5em plus 0.2em minus 0.1em}

% --- Tables and graphics ---
\usepackage{booktabs}
\usepackage{array}
\usepackage{enumitem}
\usepackage{xcolor}
\usepackage{graphicx}
\usepackage{tikz}
\usetikzlibrary{positioning,shapes.geometric,arrows.meta,fit,backgrounds,calc}
\usepackage{caption}
\usepackage{amsmath,amssymb}

% --- Code listings ---
\usepackage{listings}
\definecolor{codegray}{rgb}{0.96,0.96,0.96}
\definecolor{codeborder}{rgb}{0.85,0.85,0.85}
\definecolor{codecomment}{rgb}{0.30,0.55,0.30}
\definecolor{codekeyword}{rgb}{0.10,0.30,0.70}
\definecolor{codestring}{rgb}{0.65,0.15,0.10}

\lstdefinestyle{cstyle}{
    language=C,
    backgroundcolor=\color{codegray},
    basicstyle=\ttfamily\footnotesize,
    breaklines=true,
    frame=single,
    rulecolor=\color{codeborder},
    showstringspaces=false,
    keywordstyle=\color{codekeyword}\bfseries,
    commentstyle=\color{codecomment}\itshape,
    stringstyle=\color{codestring},
    aboveskip=1em,
    belowskip=1em,
}

\lstdefinestyle{shellstyle}{
    language=bash,
    backgroundcolor=\color{codegray},
    basicstyle=\ttfamily\small,
    breaklines=true,
    frame=single,
    rulecolor=\color{codeborder},
    showstringspaces=false,
    keywordstyle=\color{codekeyword}\bfseries,
    commentstyle=\color{codecomment}\itshape,
    stringstyle=\color{codestring},
    aboveskip=1em,
    belowskip=1em,
    columns=fullflexible,
}

\lstdefinestyle{textstyle}{
    backgroundcolor=\color{codegray},
    basicstyle=\ttfamily\footnotesize,
    breaklines=true,
    frame=single,
    rulecolor=\color{codeborder},
    showstringspaces=false,
    aboveskip=1em,
    belowskip=1em,
    columns=fullflexible,
}

\lstset{style=shellstyle}

% --- Definition / Idea / Important boxes ---
\usepackage[most]{tcolorbox}
\newtcolorbox{definicion}[1][]{
    colback=blue!5!white,
    colframe=blue!50!black,
    fonttitle=\bfseries,
    title=Definición,
    sharp corners=southwest,
    rounded corners=northwest,
    arc=2pt,
    boxrule=0.6pt,
    left=8pt,right=8pt,top=4pt,bottom=4pt,
    #1
}
\newtcolorbox{idea}[1][]{
    colback=orange!5!white,
    colframe=orange!60!black,
    fonttitle=\bfseries,
    title=Idea intuitiva,
    sharp corners=southwest,
    rounded corners=northwest,
    arc=2pt,
    boxrule=0.6pt,
    left=8pt,right=8pt,top=4pt,bottom=4pt,
    #1
}
\newtcolorbox{importante}[1][]{
    colback=red!5!white,
    colframe=red!50!black,
    fonttitle=\bfseries,
    title=Importante,
    sharp corners=southwest,
    rounded corners=northwest,
    arc=2pt,
    boxrule=0.6pt,
    left=8pt,right=8pt,top=4pt,bottom=4pt,
    #1
}

% --- Hyperlinks ---
\usepackage[colorlinks=true,
            linkcolor=blue!60!black,
            urlcolor=blue!50!black,
            citecolor=blue!60!black]{hyperref}

% --- Gantt color palette ---
\definecolor{gp1}{HTML}{4C72B0}
\definecolor{gp2}{HTML}{DD8452}
\definecolor{gp3}{HTML}{55A868}
\definecolor{gp4}{HTML}{8172B2}

% --- Title ---
\title{
    \vspace{-2.5em}
    \textbf{Resumen teórico --- Sistemas Operativos} \\[0.3em]
    \large Notas del curso 7: \emph{Planificación de uso de CPU}
}
\author{
    Tomás Rodeghiero \\
    \small Sistemas Operativos --- Universidad Nacional de Río Cuarto
}
\date{2026}

\begin{document}

\maketitle

\begin{abstract}
\noindent
Este documento es un resumen teórico del capítulo 7 de las
\emph{Notas del curso} de Sistemas Operativos (Marcelo Arroyo, UNRC),
dedicado a la planificación de uso de CPU. Se presentan las funciones
del scheduler, los criterios clásicos de planificación y los algoritmos
que los implementan: \emph{First-Come First-Served}, \emph{Shortest Job
First} (con su variante preemptiva \emph{SRTF}), \emph{Round Robin},
planificación por prioridades estáticas y dinámicas, \emph{Multilevel
Queues} y su variante con feedback (MLF). Se discuten las propiedades
de cada uno (preemption, fairness, starvation, contención), las
métricas de evaluación (turnaround, espera, respuesta, throughput) y
se incluye el caso de estudio del scheduler de Linux. La última parte
trata la planificación en multiprocesadores: arquitecturas AMP y SMP,
multi-core, hyperthreading y la tensión entre \emph{processor affinity}
y \emph{load balancing}. El objetivo es servir como material de
estudio para preparar el Práctico~4 y el segundo parcial.
\end{abstract}

\tableofcontents
\bigskip

% =====================================================================
\section{Introducción}
% =====================================================================

Una computadora moderna corre decenas de procesos a la vez con un
número limitado de CPUs. Para que todos avancen, alguien tiene que
decidir, instante a instante, \emph{quién} usa la CPU. Esa decisión es
trabajo del \textbf{scheduler} o \emph{planificador}. Las Notas del
curso (capítulo 7) lo definen así:

\begin{definicion}
El \textbf{scheduler de CPU} de un sistema operativo es el encargado
de determinar:
\begin{itemize}[leftmargin=*,nosep]
    \item \emph{cuándo} una tarea debe continuar con su ejecución
        (asignarle una CPU),
    \item \emph{dónde} (en qué procesador la ejecuta), y
    \item \emph{por cuánto tiempo}.
\end{itemize}
\end{definicion}

Estas tres preguntas no se contestan igual en todos los sistemas. En
un \emph{batch system} de los años 70 importa terminar la mayor
cantidad de \emph{jobs} por hora; en una computadora interactiva
moderna importa que la respuesta del editor de texto sea instantánea.
Por eso existen distintos algoritmos: cada uno optimiza un criterio
distinto, y el rol del diseñador del SO es elegir el que mejor sirve a
la carga de trabajo esperada.

% =====================================================================
\section{Tipos de schedulers}
% =====================================================================

Las Notas distinguen dos clases de planificadores que operan en
escalas de tiempo distintas:

\subsection{Long-term scheduler}

Aparece en \emph{batch systems} y en clusters/HPC. Los usuarios
\emph{submiten jobs} a una cola (\emph{jobs spool}), donde cada job
declara los recursos que necesita (número de CPUs, memoria, tiempo
estimado de ejecución). El long-term scheduler elige qué jobs se
admiten al sistema y en qué orden, con el objetivo típico de
\emph{minimizar el tiempo medio de espera}. El algoritmo
\textbf{Shortest Job First} (SJF) es el ejemplo paradigmático: si los
tiempos de ejecución se conocen, ejecutar primero los más cortos
minimiza la espera promedio.

\subsection{Short-term scheduler}

Es el que usan los SO modernos \emph{time-sharing}. Decide en cada
momento, entre los procesos en estado \texttt{READY}, cuál recibe la
CPU a continuación. Trabaja en escalas de milisegundos o
microsegundos. Es el scheduler que se estudia en este capítulo y el
que aparece en todos los ejercicios del Práctico~4.

\begin{idea}
La diferencia clave: el long-term scheduler decide \emph{si} un job
entra al sistema (a veces queda esperando horas en la cola); el
short-term scheduler decide qué proceso \emph{ya admitido} corre en
la CPU \emph{ahora}.
\end{idea}

% =====================================================================
\section{Criterios de planificación}
% =====================================================================

¿Qué significa que un scheduler sea ``bueno''? Las Notas listan los
criterios habituales:

\begin{itemize}[leftmargin=*]
    \item \textbf{Maximizar la utilización de la CPU.} Que la CPU esté
        ocupada haciendo trabajo útil el mayor porcentaje de tiempo
        posible. Lo ideal es 100\%; en la práctica, 40--90\% según la
        carga.
    \item \textbf{Justicia (\emph{fairness}).} Todas las tareas
        deberían recibir tiempos parejos de uso de CPU.
    \item \textbf{Throughput (rendimiento).} Número de procesos
        completados por unidad de tiempo.
    \item \textbf{Tiempo de retorno (\emph{turnaround}).} Tiempo total
        que un proceso pasa en el sistema, desde que arribó hasta que
        terminó. Suma los tiempos en estados \texttt{READY},
        \texttt{RUNNING} y \texttt{WAITING}.
    \item \textbf{Tiempo de espera.} Tiempo que el proceso pasó en
        \texttt{READY} esperando por la CPU.
    \item \textbf{Tiempo de respuesta.} Latencia desde que llega un
        requerimiento hasta que el sistema empieza a procesarlo.
        Crítico en sistemas interactivos.
\end{itemize}

Adicionalmente, un scheduler debe ser \textbf{eficiente} (mínima
sobrecarga: tomar la decisión rápido, ya que se ejecuta muchas veces
por segundo) y \textbf{general} (funcionar razonablemente en cualquier
arquitectura para distintas cargas de trabajo).

\begin{importante}
Estos criterios \textbf{se contraponen entre sí}. Maximizar throughput
puede afectar negativamente al tiempo de respuesta; favorecer fairness
puede empeorar el throughput total. Por eso no existe ``el'' mejor
algoritmo: cada uno hace bien una métrica a costa de otras.
\end{importante}

\subsection{Métricas en una corrida concreta}

Las fórmulas que se usan en los ejercicios:

\begin{align*}
\text{Espera}_i      &= \text{Inicio}_i - \text{Arribo}_i \quad (\text{algoritmos no preemptivos sin I/O}) \\
\text{Turnaround}_i  &= \text{Fin}_i - \text{Arribo}_i \\
\text{Respuesta}_i   &= \text{PrimerInicio}_i - \text{Arribo}_i
\end{align*}

En la versión preemptiva con I/O,
$\text{Espera}_i = \text{Turnaround}_i - \text{CPU}_i - \text{WAIT}_i$.

Los promedios se calculan dividiendo por la cantidad de procesos:
$\overline{W} = \frac{1}{N} \sum_i W_i$ y análogamente para
$\overline{T}$ y $\overline{R}$.

% =====================================================================
\section{Preemption: apropiativo vs.\ no apropiativo}
% =====================================================================

Una propiedad transversal a todos los algoritmos es si pueden o no
quitarle la CPU a un proceso en plena ejecución.

\begin{definicion}
Un algoritmo de planificación es \textbf{preemptivo} (apropiativo) si
el scheduler puede desalojar a un proceso que está corriendo, antes
de que termine su ráfaga de CPU. Es \textbf{no preemptivo} si, una
vez que un proceso recibe la CPU, la conserva hasta terminar la
ráfaga (o bloquearse voluntariamente en una syscall).
\end{definicion}

La preemption se implementa con la \emph{interrupción del timer}: el
hardware genera ticks periódicos, el handler del kernel decide si el
quantum se acabó y, si es así, invoca \texttt{yield()} para hacer un
context switch. Esto se conoce como \emph{preemptive multitasking} y
es la base de los sistemas time-sharing modernos.

Una alternativa histórica es el \textbf{cooperative multitasking}: los
procesos liberan la CPU \emph{voluntariamente} mediante una syscall.
Si un proceso se ``cuelga'' o no coopera, monopoliza la CPU. Por eso
los SO modernos usan preemption por timer.

% =====================================================================
\section{Algoritmos de planificación}
% =====================================================================

A partir de aquí se asume que el scheduler decide entre los procesos
\texttt{READY} (también llamados \texttt{RUNNABLE}), es decir, los
que ya están en condiciones de usar la CPU.

% ---------------------------------------------------------------------
\subsection{First-Come First-Served (FCFS)}
% ---------------------------------------------------------------------

\begin{definicion}
\textbf{FCFS} (\emph{primero que llega, primero servido}, también
llamado FIFO) es el algoritmo más simple. Se mantiene una cola FIFO
de \texttt{READY}. Cuando un proceso pasa a \texttt{READY} se encola
al final; el primero de la cola toma la CPU. Es \textbf{no
preemptivo}: el proceso conserva la CPU hasta terminar su ráfaga.
\end{definicion}

\paragraph{Ejemplo.} Supongamos los procesos
\begin{center}\small
\begin{tabular}{@{}lcc@{}}\toprule
Proceso & Arribo & Ráfaga \\\midrule
P1 & 0 & 8 \\
P2 & 0.5 & 3 \\
P3 & 3 & 5 \\\bottomrule
\end{tabular}
\end{center}

El orden de ejecución es P1, P2, P3 (orden de arribo). Diagrama de
Gantt:

\begin{center}
\begin{tikzpicture}[xscale=0.5,yscale=0.7]
    \draw[fill=gp1] (0,0) rectangle (8,1)  node[midway,white,font=\bfseries]{P1};
    \draw[fill=gp2] (8,0) rectangle (11,1) node[midway,white,font=\bfseries]{P2};
    \draw[fill=gp3] (11,0) rectangle (16,1) node[midway,white,font=\bfseries]{P3};
    \foreach \x in {0,8,11,16} \draw (\x,0) -- (\x,-0.15) node[below,font=\small]{\x};
\end{tikzpicture}
\end{center}

Esperas: $W_{P1} = 0$, $W_{P2} = 8 - 0.5 = 7.5$, $W_{P3} = 11 - 3 = 8$.
Promedio $\overline{W} = 5.17$.

\paragraph{Ventajas.}
\begin{itemize}[leftmargin=*,nosep]
    \item Simplísimo: la cola FIFO basta para implementarlo.
    \item No requiere información sobre las ráfagas.
    \item No tiene problemas de \emph{starvation} (todos avanzan en
        el orden de llegada).
\end{itemize}

\paragraph{Desventajas.}
\begin{itemize}[leftmargin=*,nosep]
    \item Suele producir \textbf{tiempos de espera promedio grandes}.
    \item Sufre el \textbf{convoy effect}: si llega primero un proceso
        muy largo, arrastra la espera de todos los que llegan después.
    \item Inutilizable en sistemas time-sharing, porque al ser no
        preemptivo, un proceso CPU-bound puede acaparar la CPU
        durante mucho tiempo.
\end{itemize}

% ---------------------------------------------------------------------
\subsection{Shortest Job First (SJF) y Shortest Remaining Time First (SRTF)}
% ---------------------------------------------------------------------

\begin{definicion}
\textbf{SJF}: el scheduler elige siempre el proceso con la
\emph{ráfaga más corta} entre los \texttt{READY}. La versión clásica
es \textbf{no preemptiva} (típica de batch systems). La versión
preemptiva, llamada \textbf{SRTF} (Shortest Remaining Time First), sí
desaloja al proceso en ejecución si llega uno con ráfaga restante
menor.
\end{definicion}

SJF se usa en \emph{long-term schedulers} y en batch systems donde el
usuario declara el tiempo estimado del job. Aplicado como
\emph{short-term scheduler}, lo correcto es llamarlo \emph{shortest
CPU-burst first}, porque lo que se considera no es la duración total
del proceso sino la duración de su \emph{próxima ráfaga de CPU} (el
intervalo entre dos bloqueos por I/O).

\paragraph{Ráfaga.} Una \emph{ráfaga (burst)} de CPU es el tiempo que
un proceso usa la CPU desde que es planificado hasta que la libera
(por bloqueo en I/O o por terminar). El comportamiento típico de un
proceso alterna ráfagas de CPU y períodos de WAIT:

\begin{lstlisting}[style=textstyle]
burst1   sleeping     burst2     sleeping  ...
       syscall  wakeup      syscall
time -->
\end{lstlisting}

\paragraph{Estimación de la próxima ráfaga.} Como SJF requiere conocer
de antemano la duración, en un short-term scheduler hay que
\emph{estimarla}. La técnica clásica es el \textbf{promedio
exponencial}: si $t_n$ es la duración real de la última ráfaga y
$\tau_n$ la última estimación, entonces
\[
\tau_{n+1} \;=\; \alpha \, t_n \;+\; (1 - \alpha) \, \tau_n,
\quad 0 \le \alpha \le 1.
\]
$\alpha$ controla el peso relativo entre la \emph{historia reciente} y
la \emph{historia pasada}. Casos extremos:
\begin{itemize}[leftmargin=*,nosep]
    \item $\alpha = 0$: $\tau_{n+1} = \tau_n$. La historia reciente
        no se tiene en cuenta.
    \item $\alpha = 1$: $\tau_{n+1} = t_n$. Solo importa la última
        ráfaga.
    \item $\alpha = 0.5$ es un balance común entre las dos.
\end{itemize}

\paragraph{Ejemplo.} Mismos procesos del FCFS pero asumiendo que
\emph{todos arriban en $t = 0$} (caso clásico para mostrar la
optimalidad de SJF). Orden por ráfaga: P2 (3), P3 (5), P1 (8).

\begin{center}
\begin{tikzpicture}[xscale=0.5,yscale=0.7]
    \draw[fill=gp2] (0,0) rectangle (3,1) node[midway,white,font=\bfseries]{P2};
    \draw[fill=gp3] (3,0) rectangle (8,1) node[midway,white,font=\bfseries]{P3};
    \draw[fill=gp1] (8,0) rectangle (16,1) node[midway,white,font=\bfseries]{P1};
    \foreach \x in {0,3,8,16} \draw (\x,0) -- (\x,-0.15) node[below,font=\small]{\x};
\end{tikzpicture}
\end{center}

Promedio de espera: $(0 + 3 + 8)/3 \approx 3.67$. Bastante mejor que
los 5.17 de FCFS sobre los mismos procesos.

\paragraph{Ventajas.}
\begin{itemize}[leftmargin=*,nosep]
    \item Es \textbf{óptimo en tiempo de espera promedio} cuando se
        conocen las ráfagas: ningún otro algoritmo logra menos espera
        promedio sobre el mismo conjunto.
    \item Reduce el \emph{convoy effect}: un proceso largo no
        arrastra a los cortos.
\end{itemize}

\paragraph{Desventajas.}
\begin{itemize}[leftmargin=*,nosep]
    \item \textbf{No se conoce} la duración de la próxima ráfaga (hay
        que estimarla).
    \item Puede producir \textbf{starvation}: si llegan continuamente
        procesos cortos, los largos pueden no ejecutarse nunca.
    \item En la versión preemptiva (SRTF), cada llegada implica una
        evaluación; la sobrecarga aumenta.
\end{itemize}

% ---------------------------------------------------------------------
\subsection{Round Robin (RR)}
% ---------------------------------------------------------------------

\begin{definicion}
\textbf{Round Robin} es la variante preemptiva de FCFS, pensada para
sistemas time-sharing. Se mantiene una cola FIFO de \texttt{READY}; al
otorgarle la CPU a un proceso se le asigna un \emph{quantum} (también
llamado \emph{time slice}). Si el proceso consume todo el quantum sin
liberar la CPU, el timer lo desaloja, vuelve a \texttt{READY} y se
encola al final. Si la libera antes (por I/O o por terminar), el
quantum no le alcanza pero igual cede la CPU al siguiente.
\end{definicion}

\paragraph{Quantum.} Es el parámetro central de RR. Hay un compromiso
clásico:
\begin{itemize}[leftmargin=*,nosep]
    \item Quantum \textbf{muy grande}: RR converge a FCFS.
    \item Quantum \textbf{muy chico}: el sistema gasta mucho tiempo en
        cambios de contexto en lugar de hacer trabajo útil. En el
        límite, aproxima un \emph{processor sharing} ideal pero con
        overhead inaceptable.
\end{itemize}
La elección habitual está entre 10 y 100 ms, y depende del costo del
context switch en la arquitectura.

\paragraph{Ejemplo.} Mismos procesos que antes, con $q = 2$:

\begin{center}
\begin{tikzpicture}[xscale=0.42,yscale=0.7]
    \draw[fill=gp1] (0,0)  rectangle (2,1)  node[midway,white,font=\bfseries\small]{P1};
    \draw[fill=gp2] (2,0)  rectangle (4,1)  node[midway,white,font=\bfseries\small]{P2};
    \draw[fill=gp1] (4,0)  rectangle (6,1)  node[midway,white,font=\bfseries\small]{P1};
    \draw[fill=gp3] (6,0)  rectangle (8,1)  node[midway,white,font=\bfseries\small]{P3};
    \draw[fill=gp2] (8,0)  rectangle (9,1)  node[midway,white,font=\bfseries\small]{P2};
    \draw[fill=gp1] (9,0)  rectangle (11,1) node[midway,white,font=\bfseries\small]{P1};
    \draw[fill=gp3] (11,0) rectangle (13,1) node[midway,white,font=\bfseries\small]{P3};
    \draw[fill=gp1] (13,0) rectangle (15,1) node[midway,white,font=\bfseries\small]{P1};
    \draw[fill=gp3] (15,0) rectangle (16,1) node[midway,white,font=\bfseries\small]{P3};
    \foreach \x in {0,2,4,6,8,9,11,13,15,16}
        \draw (\x,0) -- (\x,-0.15) node[below,font=\scriptsize]{\x};
\end{tikzpicture}
\end{center}

P2 (la ráfaga más corta) termina antes que en FCFS, pero P1 (la más
larga) termina mucho más tarde porque ejecuta en cuatro porciones
intercaladas.

\paragraph{Ventajas.}
\begin{itemize}[leftmargin=*,nosep]
    \item \textbf{Acota el tiempo de respuesta}: ningún proceso tiene
        que esperar más que $(N-1) \cdot q$ unidades antes de
        ejecutar, donde $N$ es la cantidad de procesos en
        \texttt{READY}.
    \item \textbf{Justo}: todos los procesos progresan
        regularmente.
    \item Sin starvation.
\end{itemize}

\paragraph{Desventajas.}
\begin{itemize}[leftmargin=*,nosep]
    \item \textbf{Mayor turnaround promedio} que FCFS o SJF cuando hay
        procesos largos, porque la ejecución se fragmenta.
    \item Costo de los \emph{context switches}: cada quantum agotado
        produce uno.
\end{itemize}

% ---------------------------------------------------------------------
\subsection{Planificación por prioridades}
% ---------------------------------------------------------------------

\begin{definicion}
A cada proceso se le asocia un \textbf{número de prioridad}; el
scheduler elige siempre el de mayor prioridad. La relación de orden
puede ser arbitraria: en algunos sistemas un menor valor implica
mayor prioridad (Linux real-time, 0--99); en otros, lo contrario
(Linux normal vía \texttt{nice}, $-20$ a $+19$).
\end{definicion}

\paragraph{Estática vs.\ dinámica.}
\begin{itemize}[leftmargin=*,nosep]
    \item \textbf{Estática}: la prioridad se asigna en tiempo de
        creación y no cambia. Es lo que usa el Práctico~4 Ejercicio~4.
    \item \textbf{Dinámica}: el sistema recalcula la prioridad
        periódicamente según el comportamiento del proceso. Es la
        base del scheduler 4.3BSD y, conceptualmente, de MLF.
\end{itemize}

\paragraph{Preemptiva vs.\ no preemptiva.} El algoritmo puede
implementarse en cualquiera de las dos formas: preemptivo si llegar
una tarea de mayor prioridad desaloja a la actual; no preemptivo si
la decisión solo se toma cuando la actual termina o se bloquea.

\paragraph{Ejemplo de prioridad dinámica (4.3BSD).} La fórmula
clásica es
\[
\text{priority} = \frac{\text{recent\_cpu\_usage}()}{2} + 60,
\]
donde mayor valor numérico implica menor prioridad efectiva. Cada
segundo se recalcula. Un proceso CPU-bound acumula
\texttt{recent\_cpu} y su prioridad numérica sube (es decir, baja en
términos efectivos). Un proceso I/O-bound, que se bloquea seguido,
mantiene \texttt{recent\_cpu} chico y conserva alta prioridad.
Resultado: el sistema favorece a los interactivos sin negarles
servicio a los CPU-bound, porque la penalización solo refleja la
historia reciente.

\paragraph{Starvation y aging.} El gran problema de las prioridades
es la \textbf{starvation}: un proceso de baja prioridad podría no
ejecutarse nunca si llegan continuamente otros más prioritarios. Es,
por construcción, un algoritmo \emph{no justo}.

\begin{importante}
\textbf{Aging (envejecimiento)}: incrementar gradualmente la
prioridad de los procesos que llevan mucho tiempo esperando. Tarde o
temprano todos llegan al frente de la cola, evitando starvation. Es
la solución estándar al problema y aparece explícitamente en la
teoría.
\end{importante}

% ---------------------------------------------------------------------
\subsection{Multilevel Queues (MLQ)}
% ---------------------------------------------------------------------

\begin{definicion}
En \textbf{Multilevel Queues} los procesos se clasifican en grupos.
A cada grupo se le asigna una prioridad: es entonces un algoritmo
basado en prioridades. El scheduler selecciona siempre el primer
proceso del grupo de mayor prioridad. Cada nivel puede usar un
algoritmo distinto internamente (típicamente RR).
\end{definicion}

Ejemplo común: dos grupos, ``interactivos'' (alta prioridad) y
``background / CPU-bound'' (baja prioridad). Los interactivos se
priorizan para minimizar el tiempo de respuesta.

\begin{center}
\begin{tikzpicture}[
    every node/.style={font=\ttfamily\small, draw, minimum width=1.6cm, minimum height=0.55cm},
    edge/.style={->,>=Stealth}
]
    \node (l1) at (0,2.4) {level1};
    \node[draw=none,right=0.2cm of l1.east] {$\to$ //};
    \node (l2) at (0,1.6) {level2};
    \node (p5) [right=0.5cm of l2,fill=blue!10] {p5};
    \node (p2) [right=0.3cm of p5,fill=blue!10] {p2};
    \node[draw=none,right=0.2cm of p2.east] {$\to$ //};
    \node (l3) at (0,0.8) {level3};
    \node (p3) [right=0.5cm of l3,fill=blue!10] {p3};
    \node[draw=none,right=0.2cm of p3.east] {$\to$ //};
    \node (lN) at (0,0)   {levelN};
    \node (p1) [right=0.5cm of lN,fill=blue!10] {p1};
    \node[draw=none,right=0.2cm of p1.east] {$\to$ //};
    \node[draw=none,left=0.4cm of l1.west] (hp) {\scriptsize higher priority};
    \node[draw=none,left=0.4cm of lN.west] (lp) {\scriptsize lower priority};
\end{tikzpicture}
\end{center}

% ---------------------------------------------------------------------
\subsection{Multilevel Feedback (MLF)}
% ---------------------------------------------------------------------

\begin{definicion}
\textbf{MLF} es la variante \emph{dinámica} de MLQ: los procesos se
\textbf{mueven entre niveles} según su comportamiento.
\begin{itemize}[leftmargin=*,nosep]
    \item Si un proceso \textbf{consume todo el quantum}, se considera
        CPU-bound y se \emph{baja} un nivel.
    \item Si un proceso \textbf{no consume todo el quantum} (cede
        voluntariamente, típicamente por bloqueo en I/O), se considera
        interactivo y se \emph{sube} de nivel (o se mantiene si ya
        está en el nivel más alto).
\end{itemize}
Cada nivel suele tener su propio quantum, típicamente menor en los
niveles altos (más respuesta) y mayor en los bajos (menos overhead
para los CPU-bound).
\end{definicion}

\begin{idea}
La intuición: ``los premios y los castigos son automáticos''. El
scheduler aprende solo qué procesos son interactivos (los premia con
prioridad alta) y cuáles CPU-bound (los relega a niveles bajos
donde compiten entre ellos), sin que un humano tenga que clasificarlos
a mano. Es lo más parecido a un scheduler ``adaptativo'' clásico.
\end{idea}

MLF típicamente cumple el objetivo de \emph{minimizar el tiempo de
respuesta de los procesos interactivos sin negarle progreso a los
CPU-bound}. La penalización a los CPU-bound es ``compartirse el nivel
inferior'', no quedarse afuera.

% =====================================================================
\section{Caso de estudio: el scheduler de Linux}
% =====================================================================

Las Notas dedican una sección al scheduler de Linux como ejemplo
concreto.

\subsection{Representación de procesos y threads}

Cada proceso o thread se representa con la estructura
\texttt{task\_struct}. Un proceso multi-threaded se modela con varias
instancias: una por el thread principal, una por cada thread hijo;
todas comparten el mismo \texttt{pid}.

\subsection{Rango de prioridades}

Linux usa un scheduler basado en prioridades en el rango \textbf{0 a
139}, dividido en dos grandes clases:

\begin{itemize}[leftmargin=*]
    \item \textbf{Real-time (0--99).} Para procesos del sistema y
        kernel threads que requieren garantías de tiempo. Políticas
        \texttt{SCHED\_FIFO} (no preemptiva entre procesos del mismo
        nivel) o \texttt{SCHED\_RR} (RR con quantum fijo).
    \item \textbf{Normal (100--139).} Para procesos de usuario.
        Políticas \texttt{SCHED\_NORMAL} (RR con quantum variable),
        \texttt{SCHED\_BATCH} (procesos no interactivos) y
        \texttt{SCHED\_IDLE} (los de menor prioridad).
\end{itemize}

\paragraph{Tareas de tiempo real.} Una tarea \emph{real-time} requiere
que el sistema le asigne recursos en tiempos acotados (con
vencimientos). Ejemplo clásico: una aplicación que graba un DVD
necesita mantener el buffer de datos siempre lleno, porque el láser
no puede parar de grabar. Si el SO se atrasa, la grabación falla.

\subsection{El valor \texttt{nice}}

Para procesos normales, la prioridad se determina a partir del valor
\textbf{nice}, en el rango $-20$ (más prioritario) a $+19$ (menos
prioritario):
\[
\text{priority} = \text{nice} + 20.
\]
La syscall \texttt{nice(value)} permite ajustarlo. Conceptualmente:
\begin{itemize}[leftmargin=*,nosep]
    \item Cualquier usuario puede \emph{subir} su nice (bajar su
        prioridad): cederle prioridad a otros.
    \item Solo \texttt{root} puede \emph{bajar} su nice (subir su
        prioridad), porque permitir que un usuario suba la suya
        rompería la equidad y abriría la puerta a denegación de
        servicio.
\end{itemize}
El valor de \texttt{nice} también determina el quantum asignado al
proceso. Los procesos normales arrancan con prioridad 0.

\subsection{Algoritmos a lo largo del tiempo}

Internamente Linux fue cambiando su scheduler a lo largo de los años:
\textbf{O(1) scheduler} en los 2000, \textbf{CFS} (\emph{Completely
Fair Scheduler}) entre 2007 y 2024, y desde 2024 el algoritmo
\textbf{EEVDF} (\emph{Earliest Eligible Virtual Deadline First}). El
scheduling es \textbf{preemptivo}: si arriba una tarea con más alta
prioridad que la corriente, ésta es desalojada y la nueva toma la
CPU.

Las versiones recientes permiten integrar \emph{schedulers a medida}
mediante la \emph{scheduling class} \texttt{sched\_ext}, que admite
implementar políticas como código eBPF cargado en el kernel.

% =====================================================================
\section{Multiprocesadores}
% =====================================================================

Hasta acá se asumió implícitamente que hay \emph{una} CPU. En
arquitecturas modernas hay varias: el scheduler tiene que decidir no
solo \emph{qué} ejecutar, sino \emph{en qué procesador}.

\subsection{AMP vs.\ SMP}

\begin{definicion}
\textbf{AMP (Asymmetric Multiprocessing)}: una CPU \emph{master}
ejecuta el código del kernel y sus servicios; las demás
(\emph{workers}) ejecutan procesos de usuario. El scheduler corre en
el master, lo que simplifica los problemas de paralelismo.
\end{definicion}

\begin{definicion}
\textbf{SMP (Symmetric Multiprocessing)}: todas las CPUs ejecutan
tanto código de usuario como del kernel. Es el modelo de las
arquitecturas modernas.
\end{definicion}

En SMP hay que poner especial cuidado en problemas de \textbf{concurrencia}:
varias CPUs pueden estar ejecutando código del kernel en paralelo,
así que las estructuras compartidas requieren spinlocks, instrucciones
atómicas, barreras de memoria, etc.

\subsection{Load balancing}

\begin{definicion}
\textbf{Load balancing}: emparejar la carga de trabajo entre todas
las CPUs. Si una está sobrecargada y otra ociosa, el scheduler
\emph{migra} procesos para nivelar.
\end{definicion}

Esto requiere que el sistema pueda contabilizar la carga de cada CPU
y mover procesos a las menos cargadas. En cada context switch un
proceso puede continuar en otra CPU. En Linux hay un \emph{kernel
thread} dedicado, llamado \texttt{migration}, que ejecuta
\texttt{load\_balance} para hacer este trabajo.

\subsection{Processor affinity}

La migración tiene un costo: \textbf{la caché de la CPU anterior
queda invalidada} (o, mejor dicho, deja de servir) y la caché de la
nueva tiene que volver a llenarse. Mientras se ``calienta'' la nueva
caché, el proceso ejecuta más lento porque acumula \emph{cache miss}.

\begin{definicion}
\textbf{Processor affinity}: política del scheduler que trata de
minimizar la migración de procesos entre CPUs, para preservar la
inversión en caché.
\end{definicion}

Linux ofrece dos niveles:
\begin{itemize}[leftmargin=*,nosep]
    \item \textbf{Soft affinity} (default): el scheduler trata de no
        migrar pero puede hacerlo si compensa.
    \item \textbf{Hard affinity}: la syscall
        \texttt{sched\_setaffinity} fija el proceso a un subconjunto
        específico de CPUs y el balanceador no puede sacarlo de ahí.
\end{itemize}

\begin{importante}
\textbf{Affinity y load balancing se contraponen}. Affinity pide no
migrar; load balancing pide migrar para nivelar la carga. La
solución es un compromiso: \emph{migrar solo cuando el desbalance
proyectado supera al costo esperado de la migración}. Esto se
implementa con per-CPU runqueues, balanceadores periódicos y
umbrales de desbalance.
\end{importante}

\subsection{Multi-core e hyperthreading}

\paragraph{Multi-core.} Una arquitectura multi-core replica varias
CPUs (\emph{cores}) en un mismo chip. Cada core puede usar múltiples
\emph{pipelines} para ejecutar varios threads en paralelo y
minimizar los \emph{memory stalls} (esperas por datos no cacheados).

\paragraph{Hyperthreading.} Es la capacidad de un core de ejecutar
más de un thread en paralelo, asignándoles pipelines distintos. Una
CPU de dos cores con dos pipelines cada uno se ve como un
multiprocesador de \emph{4 CPUs lógicas}.

\paragraph{Hart (RISC-V).} En la arquitectura RISC-V, un \emph{hart}
es la abstracción de un \emph{hardware thread} o \emph{logical
processor}: una unidad lógica de procesamiento con su propio conjunto
de registros. Una implementación puede materializar varios harts en
un mismo core mediante múltiples pipelines.

% =====================================================================
\section{Comparación rápida de algoritmos}
% =====================================================================

\begin{center}
\small
\begin{tabular}{@{}lccp{3.5cm}p{3.8cm}@{}}
\toprule
Algoritmo & Preemption & Convención & Ventaja principal & Desventaja principal \\
\midrule
FCFS    & No  & Cola FIFO      & Simple, sin starvation & Convoy effect, no time-sharing \\
SJF     & No  & Menor ráfaga   & Mínimo tiempo de espera promedio (óptimo) & Requiere conocer ráfagas; starvation \\
SRTF    & Sí  & Menor remanente & Igual que SJF, reactivo a llegadas & Más context switches; starvation \\
RR      & Sí  & Cola FIFO + quantum & Acota tiempo de respuesta, fair & Mayor turnaround si hay procesos largos \\
Prioridades & Cualq. & Mayor prioridad & Refleja política del sistema & Starvation; necesita aging \\
MLQ     & Sí  & Múltiples colas + prio. fija & Separa clases de tareas & Clasificación estática \\
MLF     & Sí  & Colas con prio.\ dinámica & Premia interactivos, castiga CPU-bound & Complejo de afinar \\
\bottomrule
\end{tabular}
\end{center}

% =====================================================================
\section{Resumen final}
% =====================================================================

\begin{enumerate}[leftmargin=*]
    \item El \textbf{scheduler de CPU} decide cuándo, dónde y por
        cuánto tiempo cada tarea ejecuta. Hay schedulers de
        \emph{long-term} (admisión) y \emph{short-term} (selección
        instantánea).
    \item Los \textbf{criterios} clásicos --- utilización, fairness,
        throughput, turnaround, espera, respuesta --- \emph{se
        contraponen}. No existe el algoritmo perfecto; se elige según
        la carga de trabajo.
    \item La \textbf{preemption} es la propiedad que permite al kernel
        recuperar la CPU sin la cooperación del proceso, gracias a las
        interrupciones del timer. Es la base de los sistemas
        time-sharing.
    \item \textbf{FCFS} es simple pero sufre el \emph{convoy effect}.
        \textbf{SJF} (no preemptivo) y \textbf{SRTF} (preemptivo)
        minimizan la espera promedio si se conocen las ráfagas, pero
        admiten starvation. La duración se estima con \emph{promedio
        exponencial} $\tau_{n+1} = \alpha t_n + (1-\alpha)\tau_n$.
    \item \textbf{Round Robin} acota el tiempo de respuesta usando
        un \emph{quantum}; es la elección natural para time-sharing.
        El quantum es un compromiso: muy grande $\to$ FCFS; muy chico
        $\to$ overhead de context switches.
    \item \textbf{Prioridades} reflejan la importancia que el sistema
        asigna a cada proceso. Pueden ser \emph{estáticas} o
        \emph{dinámicas}; pueden ser \emph{preemptivas} o no. Padecen
        starvation, mitigada con \emph{aging}.
    \item \textbf{MLF} combina prioridades dinámicas con RR por nivel:
        promueve a los interactivos, releva a los CPU-bound, y es lo
        más cercano a un scheduler ``adaptativo'' clásico.
    \item En \textbf{Linux} la prioridad va de 0 a 139, con clases
        real-time (0--99) y normal (100--139). El valor \texttt{nice}
        ($-20$ a $+19$) controla la prioridad y el quantum de los
        procesos de usuario. Históricamente: O(1) $\to$ CFS $\to$
        EEVDF (2024).
    \item En \textbf{multiprocesadores SMP} el scheduler debe lograr
        \emph{load balancing} (nivelar carga) y \emph{processor
        affinity} (preservar caché). Estos dos objetivos se
        contraponen y se resuelven con un compromiso: migrar solo si
        el desbalance supera el costo esperado de la migración.
\end{enumerate}

\begin{importante}
\textbf{La idea unificadora:} no hay un mejor algoritmo de
planificación. Cada criterio que se quiere maximizar (throughput,
respuesta, fairness, afinidad de caché) se contrapone con otros, y
el rol del diseñador del SO --- y del scheduler concreto que se
elija --- es navegar esos compromisos para una carga de trabajo
particular. La planificación es, en última instancia, un ejercicio
permanente de equilibrio entre objetivos contrapuestos.
\end{importante}

\end{document}
